\subsection{Semi-norms and convex sets}

Let \(\mathbf{E}\) be a vector space over \(\mathbf{R}\) ; a mapping \(p\) of \(\mathbf{E}\) in \(\mathbf{R}\) is positively homogeneous if, for every \(\lambda \geqslant 0\) and all \(x \in \mathbf{E}\) we have

\[
p (\lambda x) = \lambda p (x).\tag{5}
\]

A positively homogeneous function p on E is convex if, and only if, it satisfies axiom \((\mathrm{SN}_{\mathrm{II}})\) of II, p.~1 for all x, y of E;

\[
p (x + y) \leqslant p (x) + p (y).\tag{6}
\]

In fact, if \(p\) is convex, then for \(x, y\) in \(E\) ,

\[
p \left(\frac {1}{2} (x + y)\right) \leqslant \frac {1}{2} p (x) + \frac {1}{2} p (y)
\]

and, by (5), this relation is equivalent to (6). Conversely, if (6) holds, then we also have for all \(\lambda\) such that \(0 < \lambda < 1\) ,

\[
p (\lambda x + (1 - \lambda) y) \leqslant p (\lambda x) + p ((1 - \lambda) y) = \lambda p (x) + (1 - \lambda) p (y)
\]

using (5).

A convex positively homogeneous function on E is called sub-linear.

If \(p\) is a sub-linear function defined on \(\mathbf{E}\) ; then, by II, § 2.8, corollary, for all \(a > 0\) , the set \(\mathrm{V}(p, a)\) (resp. \(\mathrm{W}(p, a)\) ) of points \(x \in \mathrm{E}\) for which \(p(x) \leqslant a\) (resp. \(p(x) < a\) ) is convex; further this set is absorbent, since for all \(x \in \mathrm{E}\) , there exists \(\lambda > 0\) such that \(p(\lambda x) = \lambda p(x) < a\) .

There is a partial converse of this result :

PROPOSITION 22. --- Let A be a convex set, containing 0, in the vector space E. For all \(x \in \mathbf{E}\) , put

\[
p _ {\mathrm{A}} (x) = \inf _ {\rho > 0, x \in \rho \mathrm{A}} \rho\tag{7}
\]

\((0 \leqslant p_{\mathbf{A}}(x) \leqslant \infty)\) . The function \(p_{\mathbf{A}}\) satisfies

\[
p _ {\mathrm{A}} (x + y) \leqslant p _ {\mathrm{A}} (x) + p _ {\mathrm{A}} (y), \quad p _ {\mathrm{A}} (\lambda x) = \lambda p _ {\mathrm{A}} (x)\tag{8}
\]

for all \(x, y\) in \(\mathbf{E}\) and \(\lambda > 0\) . If \(\mathrm{V}(p_{\mathbf{A}}, \alpha)\) (resp. \(\mathrm{W}(p_{\mathbf{A}}, \alpha)\) ) denotes the set of \(x \in \mathbf{E}\) for which \(p_{\mathbf{A}}(x) \leqslant \alpha\) (resp. \(p_{\mathbf{A}}(x) < \alpha\) ), then

\[
\mathrm{W} (p _ {\mathrm{A}}, 1) \subset \mathrm{A} \subset \mathrm{V} (p _ {\mathrm{A}}, 1).\tag{9}
\]

If \(\mathbf{A}\) is absorbent then \(p_{\mathbf{A}}\) is finite (therefore sublinear).

Since the relations \(x \in \rho A\) and \(\lambda x \in \lambda \rho A\) are equivalent when \(\lambda > 0\) , we have \(p_A(\lambda x) = \lambda p_A(x)\) for \(\lambda > 0\) . Let \(x, y\) be two points of \(E\) . If \(x\) (resp. \(y\) ) is not absorbed by \(A\) then \(p_A(x) = +\infty\) (resp. \(p_A(y) = +\infty\) ) and the inequality \(p_A(x + y) \leqslant p_A(x) + p_A(y)\) is obviously true. Suppose there exist \(\alpha > 0\) , \(\beta > 0\) such that \(x \in \alpha A\) , and \(y \in \beta A\) ; then \(x + y \in \alpha A + \beta A = (\alpha + \beta)A\) (II, p.~8, Remark); and thus \(p_A(x + y) \leqslant p_A(x) + p_A(y)\) . The inclusion \(A \subset V(p_A, 1)\) is clearly true. The inclusion \(W(p_A, 1) \subset A\) follows because \(A\) is convex and contains 0. Finally if \(A\) is absorbent then \(p_A\) is obviously finite.

The function \(p_{A}\) defined by (7) is called the gauge of the convex set A. If A is absorbent and symmetric, then \(p_{A}\) is a semi-norm.

PROPOSITION 23. --- Let E be a topological vector space. If A is an open convex set which contains 0, then \(p_{A}\) is finite and continuous, and \(\mathbf{A} = \mathbf{W}(p_{\mathbf{A}}, 1)\) . If A is a closed convex set containing 0, then \(p_{A}\) is lower semi-continuous and \(\mathbf{A} = \mathbf{V}(p_{\mathbf{A}}, 1)\) .

If \(\mathbf{A}\) is open and contains 0, then it is absorbent. For \(x \in \mathbf{A}\) , there exists \(\rho < 1\) such that \(x / \rho \in \mathbf{A}\) , and thus \(p_{\mathbf{A}}(x) < 1\) ; this, combined with (9) gives \(\mathbf{A} = \mathbf{W}(p_{\mathbf{A}}, 1)\) . Since the convex function \(p_{\mathbf{A}}\) is bounded above in the open set \(\mathbf{A}\) , it is continuous in E (II, p.~18, prop. 21).

Suppose A is closed and contains 0. For every \(x \in E\) with \(p_{\mathbf{A}}(x) \leqslant 1\) , we have \(x \in \rho A\) for all \(\rho > 1\) , therefore \(x \in A\) since A is closed; remembering (9), this shows that \(\mathbf{A} = \mathrm{V}(p_{\mathbf{A}}, 1)\) . For all \(\mu > 0\) , \(\mu A\) is therefore the set of \(x\) such that \(p_{\mathbf{A}}(x) \leqslant \mu\) ; as \(p_{\mathbf{A}}(x) \geqslant 0\) in E, this shows that \(p_{\mathbf{A}}\) is lower semi-continuous in E (GT, IV, § 6.2).

A positive sublinear function \(p\) over \(\mathbf{E}\) is the gauge of each convex set \(\mathbf{A}\) where \(\mathrm{W}(p,1)\subset \mathrm{A}\subset \mathrm{V}(p,1)\) .

